% Part: incompleteness
% Chapter: representability-in-q
% Section: prim-rec

\documentclass[../../../include/open-logic-section]{subfiles}

\begin{document}

\olfileid{inc}{req}{pri}
\olsection{আদিম পুনরাবৃত্তির অনুকরণ}

এখন দেখানো যায়, বিটা অপেক্ষক ব্যবহার করে নিয়মিত ন্যূনীকরণের মাধ্যমে আদিম
পুনরাবৃত্তির সংজ্ঞাকে ``অনুকরণ'' করা যায়। ধরি আমাদের কাছে $f(\vec x)$ ও
$g(\vec x, y, z)$ আছে। তাহলে আদিম পুনরাবৃত্তির মাধ্যমে $f$ ও~$g$ থেকে
সংজ্ঞায়িত অপেক্ষক~$h(\vec x,y)$ হলো
\begin{align*}
h(\vec x, 0) & =  f(\vec x) \\
h(\vec x, y+1) & =  g(\vec x, y, h(\vec x, y)).
\end{align*}
% BN-SRC-220 TARGET-CORRECTION-BEGIN
\par\smallskip\noindent\textnormal{\small\textbf{উৎস-সংশোধন (BN-SRC-220):} হিমায়িত পরিচায়ক বাক্যে অপেক্ষকটির আর্গুমেন্ট ভুল করে $h(x,\vec z)$ লেখা হয়েছে, কিন্তু অব্যবহিত পরের পুনরাবৃত্ত সমীকরণ ও গোটা প্রমাণে সেটি $h(\vec x,y)$। বাংলা পাঠে সামঞ্জস্যপূর্ণ স্বাক্ষর $h(\vec x,y)$ রাখা হয়েছে।} % BN-SRC-220 TARGET-CORRECTION-END
$h$-কে শুধু মিশ্রণ ও নিয়মিত ন্যূনীকরণ ব্যবহার করে $f$ ও~$g$ থেকে সংজ্ঞায়িত
করা যায়—এটি দেখাতে হবে; সঙ্গে মৌলিক অপেক্ষক এবং মিশ্রণ ও নিয়মিত
ন্যূনীকরণের মাধ্যমে তাদের থেকে সংজ্ঞায়িত অপেক্ষকও (যেমন~$\beta$) ব্যবহার
করা যাবে।

\begin{lem}
\ollabel{lem:prim-rec}
$h$-কে যদি আদিম পুনরাবৃত্তির মাধ্যমে $f$ ও $g$ থেকে সংজ্ঞায়িত করা যায়,
তবে $f$, $g$ এবং অপেক্ষক $\Zero$, $\Succ$, $\Proj{n}{i}$, $\Add$,
$\Mult$, $\Char{=}$ থেকে মিশ্রণ ও নিয়মিত ন্যূনীকরণ ব্যবহার করে তাকে
সংজ্ঞায়িত করা যায়।
\end{lem}

\begin{proof}
প্রথমে একটি সহায়ক অপেক্ষক $\hat h(\vec x, y)$ সংজ্ঞায়িত করি। এটি সেই
ক্ষুদ্রতম সংখ্যা~$d$ ফেরত দেয় যাতে $d$ নিচের শর্তপূরণকারী একটি অনুক্রমকে
সংকেতায়িত করে:
\begin{enumerate}
\item $(d)_0 = f(\vec x)$, এবং
\item প্রত্যেক $i < y$-এর জন্য $(d)_{i+1} = g(\vec x, i, (d)_i)$,
\end{enumerate}
যেখানে এখন $(d)_i$ হলো $\beta(d,i)$-এর সংক্ষিপ্ত রূপ। অন্যভাবে বললে,
$\hat h$ অনুক্রম $\tuple{h(\vec x, 0), h(\vec x, 1), \dots, h(\vec
x, y)}$ ফেরত দেয়। $\hat h$-কে লেখা যায়
\[
\hat h(\vec x, y) = \umin{d}{(\beta(d,0) = f(\vec x) \land \bforall{i <
  y}{\beta(d,i+1) = g(\vec x, i,\beta(d,i))})}.
\]
% BN-SRC-221 TARGET-CORRECTION-BEGIN
\par\smallskip\noindent\textnormal{\small\textbf{উৎস-সংশোধন (BN-SRC-221):} হিমায়িত প্রদর্শনে ন্যূনীকরণের ধর্মটি খোলা বন্ধনীর পরে আর বন্ধ হয়নি। বাংলা পাঠে সর্বজনীন সীমাবদ্ধ শর্তটির পরে অনুপস্থিত ডান বন্ধনীটি যোগ করা হয়েছে।} % BN-SRC-221 TARGET-CORRECTION-END
লক্ষ করো, এখানে কোনো আদিম পুনরাবৃত্তি দরকার নেই, শুধু ন্যূনীকরণ। যে অপেক্ষকের
উপর ন্যূনীকরণ করি সেটি বিটা অপেক্ষকের সহায়ক উপপাদ্য
\olref[bet]{lem:beta}-এর কারণে নিয়মিত।

কিন্তু এখন
\[
h(\vec x, y) = \beta(\hat h(\vec x, y), y),
\]
তাই শুধু মিশ্রণ ও নিয়মিত ন্যূনীকরণ ব্যবহার করে মৌলিক অপেক্ষকগুলি থেকে $h$-কে
সংজ্ঞায়িত করা যায়।
\end{proof}

\end{document}
